\section{Faithful model interpretations and theory extension}\label{sec:family}

Fix a consistent theory $T\supseteq\PA$.
Let $\mathcal M_T$ be the set of all structures for the language of
arithmetic with underlying set $\N$ that satisfy $T$. Arithmetic
operations and numeral interpretations are those of each model. Completeness and
the downward L\"owenheim--Skolem theorem imply
$\mathcal M_T\ne\varnothing$: every model of $T$ is infinite, and a
countably infinite model can be transported to the carrier $\N$.

Define the product category
$\mathcal D_T=\prod_{M\in\mathcal M_T}\Pfn$.
An object is a family of sets $A=(A_M)_M$. A morphism
$f:A\rightharpoonup B$ is a family
$(f_M:A_M\rightharpoonup B_M)_M$. Composition and the
restriction, range, restriction-product, and coproduct operations are coordinatewise.
The functors of Proposition~\ref{prop:model} combine to give
$F_T:S_T\to\mathcal D_T$, defined by $F_T(X)=(X_M)_M$ and
$F_T([e]_T)=(F_M[e]_T)_M$.

\begin{theorem}\label{thm:faithful}
The functor $F_T$ is faithful and preserves restrictions, zeros, binary
restriction products and coproducts, ranges, and finite joins of
restriction idempotents. The category
$\mathcal D_T$ is split and locally connected, satisfies weak choice,
and has inverse images preserving binary joins of restriction
idempotents. Moreover,
$F_T(1_\omega)\ne0$.
\end{theorem}
\begin{proof}
Preservation follows coordinatewise from Proposition~\ref{prop:model}.
If $[e]_T\ne[f]_T$, let $E$ denote the sentence
in \eqref{eq:uniform}. The theory $T\cup\{\neg E\}$ is consistent.
Completeness and downward L\"owenheim--Skolem give
$M\in\mathcal M_T$ in which the two graphs differ. Hence
$F_M([e]_T)\ne F_M([f]_T)$, so $F_T([e]_T)\ne F_T([f]_T)$.

A restriction idempotent in $\mathcal D_T$ is a family of partial
identities on subsets. Its splitting uses those subsets as the
coordinate sets. For an object $A=(A_M)_M$, take the partial identity
on a singleton $\{x\}\subseteq A_M$ at one coordinate $M$ and the zero
morphism elsewhere. This restriction idempotent is connected. These
restriction idempotents detect the definedness and value of every coordinate of a partial morphism. Hence
$\mathcal D_T$ is locally connected.

A partial monomorphism has injective coordinate functions: test its
defining property on objects with a singleton at one coordinate
and the empty set at every other coordinate.
The partial inverses of these injections form a section, proving weak
choice. Inverse images commute with unions coordinatewise. Finally,
$\mathcal M_T$ is nonempty and each $X_M$ is nonempty, so the
identity of $F_T(X)$ differs from zero.
\end{proof}

For consistent computably enumerable $T\supseteq\PA$, the source $S_T$
is not locally connected \cite[Corollary 4.4, pp.~657--658]{DPM}.

For an inclusion $T\subseteq T'$ of consistent theories, define
$Q_{T,T'}:S_T\to S_{T'}$ by $Q_{T,T'}X=X$ and
$Q_{T,T'}[e]_T=[e]_{T'}$.
Provability is monotone, so this is a well-defined covariant functor.
It preserves restrictions, zeros, ranges, binary restriction products
and coproducts, and satisfies $Q_{T,T'}(\tau_T)=\tau_{T'}$.
Since $\mathcal M_{T'}\subseteq\mathcal M_T$, restriction of families
defines a covariant projection
$\pi_{T,T'}:\mathcal D_T\to\mathcal D_{T'}$. Fixing the assignments
$X\mapsto X_M$ once for each arithmetic structure gives a strictly
commuting square:
\[
\begin{tikzcd}[column sep=large]
 S_T \arrow[r,"Q_{T,T'}"] \arrow[d,"F_T"'] &
 S_{T'} \arrow[d,"F_{T'}"]\\
 \mathcal D_T \arrow[r,"\pi_{T,T'}"'] & \mathcal D_{T'}.
\end{tikzcd}
\]
On a morphism $[e]_T$, both composites have, at every
$M\in\mathcal M_{T'}$, the partial function defined by interpreting
$G_e$ in $M$. For $T\subseteq T'\subseteq T''$, the quotient
functors satisfy $Q_{T',T''}Q_{T,T'}=Q_{T,T''}$, and the coordinate
projections satisfy
$\pi_{T',T''}\pi_{T,T'}=\pi_{T,T''}$.
The identity cases follow from the same definitions.

\subsection{Inseparable pairs under theory extension}

\begin{corollary}\label{cor:pair}
Let $T\supseteq\PA$ be consistent.
Fix $U=\omega$ and a Turing morphism
$\tau:U\times U\rightharpoonup U$ in $S_T$.
Choose an isomorphism $b_0:U\to U+U$ in $S_{\PA}$ and put
$b=Q_{\PA,T}(b_0)$. Define
\[
 d=\tau\Delta_U,\qquad
 a_0=(bd)^*(1_U+0_U),\qquad a_1=(bd)^*(0_U+1_U).
\]
Then $a_0,a_1$ are inseparable in $S_T$. For every consistent
extension $T'\supseteq T$, their images under $Q_{T,T'}$ form the
inseparable pair defined by the same formulas in $S_{T'}$.
\end{corollary}
\begin{proof}
Corollary~\ref{cor:model-pair} gives the first assertion, since
$\mathcal M_T\ne\varnothing$.
The functor $Q_{T,T'}$ preserves the operations defining the pair.
It is the identity on objects and surjective on each hom-set.
Consequently it preserves a Turing morphism: a preimage of a
target morphism has a Turing factorization in $S_T$, whose image is
the required factorization in $S_{T'}$.
Applying Corollary~\ref{cor:model-pair} to $T'$ gives inseparability
of the image pair.
\end{proof}

The morphism $F_T(\tau_T)$ is not a Turing morphism in $\mathcal D_T$.
Fix $M\in\mathcal M_T$. The partial functions
$x\mapsto F_M(\tau_T)(a,x)$, indexed by all $a\in|M|$, form a
countable family, whereas there are uncountably many partial
functions $|M|\rightharpoonup|M|$.
Choose one outside this family and use it at coordinate $M$,
with zero partial functions at the other coordinates.
The resulting morphism
$\mathbf1\times F_T\omega\rightharpoonup F_T\omega$,
where $\mathbf1$ is the family of singleton sets, cannot factor
through $F_T(\tau_T)$ with a total parameter
$\mathbf1\to F_T\omega$.

The restriction lattices of $\mathcal D_T$ are Boolean, so every image
of a disjoint source pair has a separator. Lemma~\ref{lem:switch}
supplies lifts for sections associated with complementary restriction
idempotents of $S_T$.


\subsection{Provability of numeral instances}

Suppose in addition that $T$ is computably enumerable.
For a restriction idempotent
$c=[C]_T$ on $\omega$, define the subset of the standard natural numbers
\begin{equation}\label{eq:observation}
 E_T(c)=\{n\in\N:T\vdash C(\bar n)\}.
\end{equation}
It is computably enumerable by proof enumeration and is independent of
the representative. The operation is monotone, preserves finite meets,
and preserves top and bottom. Choose a constant-function index
$\operatorname{const}(n)\in\N$ for which
$\PA\vdash\forall x\forall y\,
(G_{\operatorname{const}(n)}(x,y)\leftrightarrow y=\bar n)$, and put
$\kappa_n=[\operatorname{const}(n)]_T:\omega\to\omega$. Then
\begin{equation}\label{eq:faithful-observation}
 n\in E_T(c)
 \quad\Longleftrightarrow\quad c\kappa_n=\kappa_n
 \quad\Longleftrightarrow\quad
 F_T(c)F_T(\kappa_n)=F_T(\kappa_n).
\end{equation}
The composition formula makes
$G_{\operatorname{comp}(j_C,\operatorname{const}(n))}(x,y)$
$T$-provably equivalent to $C(\bar n)\land y=\bar n$, which gives
the first equivalence.
The second follows from faithfulness of $F_T$.

\subsection{The theory of the standard model}

For $T=\operatorname{Th}(\N)$, the identity-on-objects assignment
$D:S_{\operatorname{Th}(\N)}\to\mathrm{Rec}$, defined by
$D([e]_{\operatorname{Th}(\N)})=\varphi_e$, is an isomorphism:
two classes are equal exactly when the sentence in \eqref{eq:uniform}
is true in $\N$, equivalently when their interpreted partial functions
are equal.
The functor $F_T$ still exists, and $S_T$ has inseparable pairs,
although $T$ is complete.
Thus inseparability in $S_T$ alone does not imply incompleteness of $T$.

